\documentclass[]{../../../../tex/classes/styledArticle}
\usepackage{../../../../tex/packages/hebrewSupport}
\usepackage{../../../../tex/packages/mathShortcuts}
\usepackage{../../../../tex/packages/theoremsSupport}

\author{שחר פרץ}
\title{אלגוריתמים $\sim$ \textit{הרצאה 8}}
\begin{document}
	\maketitle
	\textbf{מרצה: }עמית ווינשטיין
	
	תזכורת: זרימה ברשתות. נתון גרף מכוון $G = (V, E)$, וקיבול $c(u, v) > 0$ על הקשתות (לא מעניינות אותנו קשתות עם קיבול $0$). יש לנו מקור $s$ ויעד $t$. נרצה להזרים כמה שיותר מ־$s$ ל־$t$. מגדירים פונקציית זרימה $f \co V \times V \to \R$, שמגדירה כמה נזרים לכל קשת. הזרימה צריכה לקיים את אילוץ הקיבול: לא נזרים יותר מהקיבול, כלומר $f(u, v) \le c(u, v)$, ישנה אנטי־סימטריה $f(u, v) = -f(v, u)$, וחוק שימור החומר $\forall v \in V \setminus \{s, t\} \co \sum_{u \in V} f(v, u) = 0$. את כל ההגדרות הללו הכללנו לקבוצות של קודקודים. 
	
	ראינו שני אלגוריתמים למצוא זרימה: 
	\begin{itemize}
		\item הגדרנו \textit{רשת שיורית} להיות כמה ניתן עוד להזרים על כל קשת. הסיבה שהאלגוריתמים החמדניים עובדים היא שאנחנו מאפשרים לאלגוריתם ''להתחרט`` ולדחוף אחורה באנטי־קשתות עם קיבול שלילי. אכן בקשת השיורית יש קשתות ואנטי־קשתות של $G$. נגדיר את הקיבול השיורי להיות $c_{f}(u, v) = c(u, v) - f(u, v)$ ב־$G_f$. 
		\item כל פעם מצאנו מסלול ב־$G_f$ והוספנו זרימה דרכו. 
		\item כל פעם שמוצאים מסלול, לפחות קשת אחת קריטית נהיית רוויה. 
	\end{itemize}
	קיבלנו חסם על איטרציות של $\oc(\sof E \cdot \sof V)$ אם בוחרים מסלול קצר ביותר (בין שתי פעמים ש־$(u, v)$ קריטית המרחק מ־$s$ ל־$u$ קטן ב־$2$, וזה מונוטוני). 
	
	כדי להראות שהתוצאה היא אכן מקסימלית, הוכחנו את משפט max-flow-min-cut. בהקשר של זרימה, חתך הוא חלוקה של הקודקודים לשתי קשתות, וסופרים כמה קיבול יש במערך (ברור שכל חתך הוא חסם עליון על כל זרימה). המשפט בא ואומר שערך זרימה גדול ביותר שווה לערך חתך קטן ביותר. 
	
	האלגוריתם של Edmonds Karp פתר את הבעיה ב־$\oc(\sof E^2 \sof V)$. שאלה מתבקשת היא האם אפשר בכל איטרציה לנצל את העובדה שהמרחק מ־$s$ ל־$t$ עולה ובכך לא לבנות כל פעם קשת שיורית מחדש. 
	
	לשם כך נגדיר את $L_f$, ה‏־layered residual network או משהו כזה. טוב וויתרתי על לסכם
	
	\subsection*{אלגוריתם Goldberg Tarjan}
	חזרתי על לסכם כי אני רואה שהחומר הזה לא בסיכום של הקורס. האלגוריתם עובד עם preflow שבשונה מ־flow מרשה שיהיה exess בקודקודים שונים מ־$s, t$. הדרישה של אילוץ הקיבול עדיין מתקיימת, ונגדיר את הקודקודים ה''פעילים`` להיות אלו שיש בהם excess חיובי, כלומר: 
	\[ \sum_{u \in V}f(u, v) - f(v, u) > 0 \]
	יש לנו צעדי push – לדחוף זרימה מגובה $h(u)$ ל־$h(v) = h(u) - 1$ על קשת $(u, v)$. לא תמיד נעשה push שירווה את כל הקשת. באלגוריתם נזרים את המינימום מבין כל ה־excess והקיבול השיורי של הקשת $(u, v)$. זה ישרה לנו את ההגדרות הבאות: 
	\begin{itemize}
		\item saturating push – עם הזרמנו את כל $c_f(u, v)$
		\item non-saturating push – אחרת (הזרמנו לפי ה־excess ככל יכולתנו)
		\item relevel – אם עבור קודקוד פעיל לא הצלנו לדחוף את כל היתרה שלו, ''נגביה`` אותו לגובה המינימלית שבנינו ברשת השיורית פלוס $1$, ע''מ שנוכל לדחוף. 
	\end{itemize}
	במידה ועשינו non-saturating push, הקודקוד לא יהיה פעיל יותר (כי הקיבול השיורי שלו יהפוך להיות $0$). 
	
	כאשר שואלים את עצמנו איך הופכים את זה לאלגוריתם, כדאי לשאול את עצמנו מאיזה preflow התחלתי מתחילים. לאחר מכן נשאל את עצמנו איזה שמורות האלגוריתם מקיים כדי שזה אשכרה יעבוד. 
	
	צריכים level התחלתי לכל קודקוד, וגם preflow התחלתי. נתחיל עם: 
	\[ f(u, v) = \begin{cases}
		c(u, v) & u = s \\
		0 & \other
	\end{cases} \]
	כלומר, מזרימים את כל מה שאפשר מ־$s$. גובה $h$ התחלתי לקודקודים יהיה: 
	\[ h(u) = \begin{cases}
		n & u = s \\
		0 & \other
	\end{cases} \]
	(כאשר כמובן $n = \sof V$). נדרוש מ־$h$, שלכל קשת שיורית $(u, v)$ מתקיים $h(u) \le h(v) + 1$. כמובן ש־$h(s) = \sof V = n$ ו־$h(t) = 0$. בהתחלה אין קשתות שיוריות שנכנסות ל־$s$ (כי הזרמנו מ־$s$ את כל מה שהיה בו). לכל אורך הריצה, אין מסלול משפר $s \squig t$ ברשת השיורית (אחרת, מהשמורה, מסלול באורך $k$ היה גורם לכך ש־$h(s) = h(t) + k$ אך $h(s) < h(t)$). באמצעות טיעון דומה לכל $v$ יתקיים $h(v) <2n$ (אם הוא יטפס מעבר לכך, יש מסלול באורך יותר מ־$n$ שיורד ל־$s$). 
	
	טענה נוספת: העדכונים משמרים את תכונות $h$. 
	\begin{itemize}
		\item ב־push, הגובה לא משתנה, ולכן הוא לא פוגע ב־$h$. רק צריך לוודא שמה שקורה בקשתות $(v, u)$ (אחרי הזרמה מ־$u$ ל־$v$ דרך $(u, v)$) שנהפכה. אכן $h(v) - h(u) = 1$ וסיימנו. 
		\item ב־relabel משפרים את $u$ בדיוק ל־$\min h(v) + 1$ ולכן מתקיים. 
	\end{itemize}
	נעצור כשאין קודקודים ''פעילים`` ובעצם אין יתרה בקןדקןדים שהם לא $s, t$ ולכן $f$ זרימה תקינה. בנוסף, אין מסלול משפר, ולכן $f$ מקסימלית. 
	
	בנוגע לזמן ריצה, נראה טענה לפיה הגובה המקסימלי לא יעבור $2 \sof V - 1$. ההוכחה פשוט. אם קודקוד פעיל יש מסלול ממנו חזרה ל־$s$, מסלול תמיד באורך לכל היותר $\sof V - 1$, וגובה $s$ הוא $\sof V$. ולכן נקבל $h(u) \le h(2) -+ \sof V - 1 = 2 \sof V - 1$. 
	
	\begin{itemize}
		\item יש לנו $\sof V$ קודקודים שנעלה לכל היותר $2 \sof V - 1$ פעמים. סה''כ לכל היותר $\oc(\sof V^{2})$ פעמים נעשה relabel. 
		\item כמה saturating push יש לנו? בין שתי פעולות כאלה על אותה הקשת, בהכרח השתמשנו בקשת ההפוכה ולכן הגובה שלה עלה ב־$2$. יש לכל היותר $\oc(\sof V)$ כאלו לכל קשת, סה''כ$\oc(\sof E \cdot \sof V)$. 
		\item כמה non-saturating pushes? נגדיר פונקציית פוטנציאל $\Phi = \sum_{v \in \mathrm{Active}}h(v)$. היא כמובן תמיד אי־שלילית. 
		\begin{itemize}
			\item relabel יכול להגדיל (למעשה הוא תמיד מגדיל) את $\Phi$, אבל סה''כ $\oc(\sof V^{2})$. 
			\item saturating push יכול להפוך קודקוד לפעיל $\sof E \sof V$ פעמים להגדיל $\sof V$ פעמים. סה''כ $\oc(\sof E \cdot \sof V^2)$
			\item non-saturating pushesהופכים את $u$ ללא פעיל, ו־$v$ בטוח יהיה פעיל (אולי עוד ממקודם). מכיוון ש־$h(v) = h(u) - 1$ בהכרח הקטנו את $\phi$ בלפחות $1$. 
		\end{itemize}
		סה''כ יש $\oc(\sof E \sof V^{2})$ פעולות. 
	\end{itemize}
	
	שתי הערות צד: 
	
	\begin{itemize}
		\item אם המרחק מ־$u$ ל־$t$ הוא $k$, בכל מקרה בסוף האלגוריתם יהיה $h(u) \ge k$. לפעמים מועיל להעלות את הגבהים באמצעות BFS אחת לכמה זמן כדי לשפר ביצועים. 
		\item אם בוחרים בכל צעד את הקודקוד הפעיל הגבוה ביותר, הסיבוכיות יורדת ל־$\oc(\sof V^2 \sqrt{E})$. 
	\end{itemize}
	
	
	
	
	\ndoc
\end{document}